\section{Why the Sentence Has No End---and Why the Body Rises}

No aspect of hell tests the human mind more sharply than its eternity. How can a mortal sin committed in an hour, or even across a lifetime of seventy years, end in an unending sentence? And why, if the soul suffers immediately after death, must the body also be raised from the dust to share in eternal punishment?

Both questions receive their first answer from the plain words of Christ Himself. In Matthew 25:46, our Lord sums up the Last Judgment in a single, balanced sentence: ``And these shall go into everlasting punishment: but the just, into life everlasting''---\textit{Et ibunt hi in supplicium aeternum: justi autem in vitam aeternam}. In Greek, the exact same adjective, \textit{aionios}, modifies both punishment and life; in Latin, the exact same word, \textit{aeternum}, governs both clauses. In \textit{The City of God} XXI.23, St.\ Augustine lays bare the inescapable logic of that single verse:

\begin{quote}
Then what a fond fancy is it to suppose that eternal punishment means long-continued punishment, while eternal life means life without end, since Christ in the very same passage spoke of both in similar terms in one and the same sentence, ``These shall go away into eternal punishment, but the righteous into life eternal!'' If both destinies are ``eternal,'' then we must either understand both as long-continued but at last terminating, or both as endless. For they are correlative,---on the one hand, punishment eternal, on the other hand, life eternal. And to say in one and the same sense, life eternal shall be endless, punishment eternal shall come to an end, is the height of absurdity. Wherefore, as the eternal life of the saints shall be endless, so too the eternal punishment of those who are doomed to it shall have no end.\endnote{St.\ Augustine, \textit{The City of God} XXI.23, trans.\ Dods.}
\end{quote}

Augustine had already exposed the same fracture six chapters earlier when answering what he calls the ``tender-hearted Christians'' (\textit{misericordes}) who hoped that all condemned men, though not the demons, would eventually be released after a fixed term of fire. If mercy alone dissolves the Lord's sentence for men, Augustine asks, why should that stream of pity dry up before it reaches the fallen angels? Origen at least followed the premise to its logical collapse by proposing the eventual restoration even of the devil---and thereby fell into a cycle of endless falls and recoveries that stripped even the saints of secure beatitude.\endnote{St.\ Augustine, \textit{The City of God} XXI.17, trans.\ Dods.} Moreover, Christ's own sentence in Matthew 25:41 binds impenitent men to the very same fire prepared for the devil and his angels. Either the divine sentence holds true for both, or human conjecture claims authority to overturn the Word of God.

Yet Catholic theology does not stop at grammar and authority; it enters into the justice of an unending sentence. Why is it righteous that a brief act of sin should incur an everlasting penalty?

Augustine answers first from the universal practice of human law (\textit{The City of God} XXI.11). No court on earth measures the length of a punishment by the seconds needed to commit the crime. A murder takes a moment, an act of treason or adultery a brief space of time, yet human justice visits such offenses with years of imprisonment, perpetual banishment, or capital execution---whose penalty lies not in the fleeting instant of the axe, ``but in this, that the offender is eternally banished from the society of the living.''\endnote{St.\ Augustine, \textit{The City of God} XXI.11, trans.\ Dods.} Human punishments fail to be eternal only because human life is brief.

St.\ Thomas Aquinas goes deeper, into the nature of moral order itself (\textit{Summa Theologiae} I-II, q.\ 87, a.\ 3). Every sin incurs punishment because it disrupts an order, and the debt of punishment endures as long as the disruption remains. Some disruptions are reparable because the governing principle of the order stays intact---as a cloud over the eye can be cured while the faculty of sight survives. Other disruptions are irreparable in themselves because they destroy the very principle of the order. In the spiritual life, the principle uniting man's will to God is charity, adhesion to the ultimate end:

\begin{quote}
Consequently if a sin destroys the principle of the order whereby man's will is subject to God, the disorder will be such as to be considered in itself, irreparable, although it is possible to repair it by the power of God. Now the principle of this order is the last end, to which man adheres by charity. Therefore whatever sins turn man away from God, so as to destroy charity, considered in themselves, incur a debt of eternal punishment.\endnote{St.\ Thomas Aquinas, \textit{Summa Theologiae} I-II, q.\ 87, a.\ 3, trans.\ Fathers of the English Dominican Province.}
\end{quote}

During earthly life, God's grace can heal that mortal wound through repentance and the sacraments; after death, the sinner remains permanently without charity, so the disorder endures forever and its penalty endures with it. In his first reply to that article, Aquinas adds a luminous sentence from St.\ Gregory the Great's \textit{Dialogues} (IV.44): ``it is just that he who has sinned against God in his own eternity should be punished in God's eternity.'' A man who fixes his final end in sin and dies unrepentant has willed sin without term, for ``the wicked would wish to live without end, that they might abide in their sins for ever.''\endnote{St.\ Thomas Aquinas, \textit{Summa Theologiae} I-II, q.\ 87, a.\ 3, ad 1, quoting St.\ Gregory the Great, \textit{Dialogues} IV.44.} And in his third reply, answering the verse from Wisdom that God does not delight in the destruction of the living (Wis 1:13), Aquinas states a principle of enduring sobriety: ``God does not delight in punishments for their own sake; but He does delight in the order of His justice, which requires them.''\endnote{St.\ Thomas Aquinas, \textit{Summa Theologiae} I-II, q.\ 87, a.\ 3, ad 3.}

For the same reason, Catholic doctrine distinguishes sharply between mortal and venial sin. In Article 5 of the same question, appealing to St.\ Gregory the Great's teaching in the \textit{Dialogues} (IV.39) that certain slighter sins are remitted after this life, Aquinas explains that not every sin destroys charity or turns the will against its final end. Where a sin introduces disorder only among temporal goods referred to the end---as when a man is overly fond of a created thing yet would not break a divine commandment for its sake---the principle of order remains intact, and ``such sins do not incur everlasting, but only temporal punishment.''\endnote{St.\ Thomas Aquinas, \textit{Summa Theologiae} I-II, q.\ 87, a.\ 5, \textit{sed contra} and corpus (citing St.\ Gregory the Great, \textit{Dialogues} IV.39), trans.\ Fathers of the English Dominican Province.} Only mortal sin, which severs charity, incurs a debt of eternal punishment.

St.\ John Chrysostom presses the same truth in pastoral terms when preaching on 1 Corinthians 3. If a single transgression of Adam in Paradise brought death upon the whole human race across millennia, how can those who have received not an earthly garden but the promise of heaven, the Blood of Christ, and daily calls to repentance complain against justice if they persist in insult to the end?\endnote{St.\ John Chrysostom, \textit{Homilies on First Corinthians} 9.2, NPNF1 vol.\ 12.} And St.\ Augustine draws a decisive proof from the liturgy of the Church itself (\textit{The City of God} XXI.24). The Church prays fervently for all her living enemies---heretics, persecutors, and hardened sinners---because while breath remains, God can still grant repentance. Yet the Church has never prayed for the devil and his angels, nor does she offer the Sacrifice for those who die manifestly outside Christ's kingdom in final impenitence; for ``the man who was not translated into Christ's kingdom while he was in the body, is now judged to be of Satan's following.''\endnote{St.\ Augustine, \textit{The City of God} XXI.24, trans.\ Dods.}

This eternity of the sentence also illuminates the Catholic dogma of the bodily resurrection of the wicked. Our Lord proclaims in the Gospel of John that ``the hour cometh wherein all that are in the graves shall hear the voice of the Son of God. And they that have done good things shall come forth unto the resurrection of life: but they that have done evil, unto the resurrection of judgment'' (John 5:28--29; cf.\ Dan 12:2). The Fourth Lateran Council defines that ``all shall rise with their own bodies which they now bear,'' and \textit{Benedictus Deus} insists that although the souls of the lost descend immediately into hell after death, ``nevertheless on the day of judgment all men will appear with their bodies before the tribunal of Christ.''\endnote{Fourth Lateran Council, \textit{Firmiter credimus}, DS 801; Benedict XII, \textit{Benedictus Deus}, DS 1002.}

Already around the year 155, St.\ Justin Martyr set this Christian faith directly over against pagan philosophy in his \textit{First Apology} to the Roman Emperor Antoninus Pius: ``And Plato, in like manner, used to say that Rhadamanthus and Minos would punish the wicked who came before them; and we say that the same thing will be done, but at the hand of Christ, and upon the wicked in the same bodies united again to their spirits which are now to undergo everlasting punishment; and not only, as Plato said, for a period of a thousand years.''\endnote{St.\ Justin Martyr, \textit{First Apology} 8, ANF vol.\ 1.}

Two centuries later, St.\ Cyril of Jerusalem instructed catechumens preparing for Easter Baptism on why the body of the sinner must rise incorruptible:

\begin{quote}
We shall be raised therefore, all with our bodies eternal, but not all with bodies alike: for if a man is righteous, he will receive a heavenly body, that he may be able worthily to hold converse with Angels; but if a man is a sinner, he shall receive an eternal body, fitted to endure the penalties of sins, that he may burn eternally in fire, nor ever be consumed. And righteously will God assign this portion to either company; for we do nothing without the body. We blaspheme with the mouth, and with the mouth we pray. With the body we commit fornication, and with the body we keep chastity. With the hand we rob, and by the hand we bestow alms; and the rest in like manner. Since then the body has been our minister in all things, it shall also share with us in the future the fruits of the past.\endnote{St.\ Cyril of Jerusalem, \textit{Catechetical Lectures} 18.19, NPNF2 vol.\ 7.}
\end{quote}

St.\ John of Damascus repeats Cyril's moral reasoning with crisp economy: if the soul alone contested in virtue, the soul alone would be crowned; if the soul alone reveled in vice, the soul alone would be punished; ``but since the soul does not pursue either virtue or vice separate from the body, both together will obtain that which is their just due.''\endnote{St.\ John of Damascus, \textit{Exposition of the Orthodox Faith} IV.27, trans.\ Salmond, NPNF2 vol.\ 9.} St.\ Alphonsus Liguori notes how our Lord's enigmatic sentence immediately following the triple refrain in Mark's Gospel---``For every one shall be salted with fire'' (Mark 9:48)---captures the terrible permanence of that risen flesh: just as salt preserves flesh from dissolving into decay, so the penal fire of hell preserves the damned in existence rather than granting them the escape of annihilation.\endnote{St.\ Alphonsus Maria de Liguori, \textit{Preparation for Death}, Consideration XXVII, First Point, p.\ 329.}

In the very next paragraph of his catechesis (18.20), Cyril draws a direct moral conclusion from the resurrection of the flesh: reverence for the body right now. ``Therefore, brethren, let us be careful of our bodies, nor misuse them as though not our own.'' The sins that scar both soul and body are washed away in Baptism and healed in Penance; the body is destined for resurrection, either to shine with the angels or to endure the fire without being consumed.
